\documentclass[11pt]{article}
\usepackage[letterpaper,margin=1in]{geometry}
\usepackage[T1]{fontenc}
\usepackage{lmodern,microtype,amsmath,amssymb,amsthm,mathtools}
\usepackage[colorlinks=true,linkcolor=blue,citecolor=blue,urlcolor=blue]{hyperref}
\numberwithin{equation}{section}
\newtheorem{theorem}{Theorem}[section]
\newtheorem{lemma}[theorem]{Lemma}
\newtheorem{proposition}[theorem]{Proposition}
\newcommand{\E}{\mathbb E}
\newcommand{\norm}[1]{\left\lVert#1\right\rVert}
\allowdisplaybreaks[1]
\setlength{\emergencystretch}{2em}
\title{Second-order bounds for graph signings}
\author{}
\date{September 29, 2026}
\begin{document}
\maketitle
\begin{abstract}
Let $\gamma_d$ be the supremum, over finite simple graphs of maximum
degree at most $d$, of the minimum signed adjacency norm. We give a
proposed proof of $\gamma_d\le 2\sqrt{d-1}+O(d^{1/6})$.
The proof combines determinant-weighted vertex insertion with a
regularized resolvent comparison, a second-moment estimate, and a
coupled scalar inequality that only requires squared lower-tail
control. We also give a proposed construction showing
$\liminf_{d\to\infty}(d-1)^{7/2}(\gamma_d-2\sqrt{d-1})\ge1/4$.
These bounds do not determine the asymptotic order of the excess.
\end{abstract}
\begin{center}\small\textit{Research draft. These proposed proofs have not
undergone independent human review.}\end{center}

\section{Statement and scope}
For a finite simple graph $G$, let
\[
 s(G)=\min_{\sigma:E(G)\to\{-1,1\}}\norm{A_\sigma(G)},
 \qquad \gamma_d=\sup_{\Delta(G)\le d}s(G).
\]
Here $\norm{\cdot}$ is the Euclidean operator norm. The supremum is
unchanged if one restricts to $d$-regular graphs, by an induced
regular completion described below. The supremum also equals the
limit of the corresponding maxima over graphs on at most $n$
vertices as $n\to\infty$. All estimates here are uniform in graph order.

\begin{theorem}\label{thm:upper}
There is an absolute constant $C$ such that, for every sufficiently
large integer $d$, every finite simple graph of maximum degree at
most $d$ has an edge signing with
\[
 \norm{A_\sigma}<2\sqrt d+C d^{1/6}.
\]
Consequently $\gamma_d\le2\sqrt{d-1}+O(d^{1/6})$.
\end{theorem}

The lower-bound construction below gives, more precisely,
\[
 \left(\frac14-o(1)\right)(d-1)^{-7/2}
 \le\gamma_d-2\sqrt{d-1}\le C d^{1/6}.
\]
There remains a substantial gap between these estimates.

The determinant weight and vertex-insertion identities used here
are inherited from the framework of Jadbabaie, Saberi, and
Sra~\cite[Section~6]{JSS}. The new proposed quantitative argument
uses the same weighted measure for all terms of a resolvent
comparison, avoiding a loss in the main contraction under vertex
deletion. A second-moment estimate then feeds back into the variance
of a fresh incident-sign quadratic form. The final insertion bound
clips its positive deviations, so no unbounded positive-tail moment
is required. No interlacing-family theorem is used.
The lower-bound construction develops Xu's triangle-and-tree
obstruction~\cite{Xu}. Neither priority nor correctness of the
extensions asserted here has been externally verified.

\section{Determinant weights, cofactors, and path trees}
Fix $R>2\sqrt d$. For each signed graph $H$ put
\[
 W_H=\det(I-A_H^2/R^2)\mathbf1_{\{\norm{A_H}<R\}},
 \qquad Z_H=\E W_H,\qquad Z_\varnothing=1,
\]
where expectations without a subscript use independent uniform edge
signs. A signing with $\norm{A_H}<R$ is called good. When $Z_H>0$,
let $\nu_H$ have density $W_H/Z_H$ relative to uniform signs. Define
\[
 \mathcal J_H(v)=\E\left[W_H\left((I-A_H/R)^{-1}_{vv}
                       +(I+A_H/R)^{-1}_{vv}\right)
                         \mathbf1_{\{\norm{A_H}<R\}}\right].
\]
The integrand is defined as zero on nongood signings; no inverse is
used there. All parameters are fixed under vertex deletion.

For $K=H-v$, fix a good core signing and let $s$ be the column of
incident signs, zero outside $N_H(v)$. Set
\[
 q_\pm=R^{-2}s^T(I\pm A_K/R)^{-1}s,\qquad
 \delta_\pm=1-q_\pm,\qquad
 \alpha=(\delta_+)_+(\delta_-)_+,
 \qquad D=(\delta_++\delta_-)\mathbf1_{\{\delta_+>0,\delta_->0\}}.
\]
Schur complementation of $I\pm A_H/R$ shows that the extension is
good exactly when both residuals are positive. It also gives
\begin{equation}\label{eq:insertion-identities}
 W_H=W_K\alpha,\qquad
 \frac{Z_H}{Z_K}=\E_*\alpha,\qquad
 \frac{\mathcal J_H(v)}{Z_K}=\E_*D,
\end{equation}
where $\E_*$ uses $\nu_K$ and independent incident signs. If the core
is nongood, every extension is nongood, since its adjacency matrix
is a principal submatrix of the extended adjacency matrix. Thus
$W_H\le W_K$, because $0\le\alpha\le1$ on good cores. In particular,
if $Z_H/Z_K\ge1/8$, the density of $\nu_H$ relative to
$\nu_K$ times independent incident signs is at most $8$.
The weight is invariant under reversal of every edge sign.

Let $\mathcal T(H,v)$ be the rooted tree of simple paths in $H$
starting at $v$, with an edge when one path extends the other by
one vertex. Its adjacency norm is at most $2\sqrt d$: orient away
from the root and sum
$2|z_pz_w|\le z_p^2/\sqrt d+\sqrt d\,z_w^2$ over edges.
For any $\zeta>2\sqrt d$, set
\[
 b_H(v)=\bigl[(I-A_{\mathcal T(H,v)}/\zeta)^{-1}\bigr]_{vv}.
\]
Removing the root gives a copy of $\mathcal T(H-v,w)$ for each
neighbor $w$. Therefore
\begin{equation}\label{eq:path-recursion}
 b_H(v)=\frac1{1-\zeta^{-2}\sum_{w\sim v}b_{H-v}(w)}.
\end{equation}
If $p=d/\zeta^2<1/4$ and $a=(1+\sqrt{1-4p})/2$, induction in
this recursion gives $1\le b_H(v)\le1/a$. Deleting a different
vertex removes a rooted subtree, so
$b_{H-u}(v)\le b_H(v)$. This follows also by termwise comparison
of the convergent nonnegative closed-walk series. These are the
only properties of path trees needed below.
\section{A quantitative cavity estimate with a second-moment bootstrap}

Fix $0<\tau\le 1/2$, $0<\beta\le1/2$, and an integer $d\ge1$.
Put
\[
 p=\frac{1-\beta^2}{4},\qquad
 \zeta^2=\frac d p,\qquad
 k=\tau\beta^2,\qquad R^2=\frac{\zeta^2}{1+k},\qquad
 a=\frac{1+\beta}{2},\qquad B=\frac1a.
\]
In particular $R>2\sqrt d$, because
$(1-\beta^2)(1+\tau\beta^2)<1$.
All determinant weights and Gibbs measures below use the radius $R$.
The deterministic reference
\[
 b_H(u)=\bigl[(I-A_{\mathcal T(H,u)}/\zeta)^{-1}\bigr]_{uu}
\]
uses the larger radius $\zeta$. Its path-tree recursion gives
$1\le b_H(u)\le B$, and deletion of any vertex $v\ne u$ gives
$b_{H-v}(u)\le b_H(u)$: the former path tree is a rooted subtree
of the latter and all its closed-walk contributions are nonnegative.

Let $\mathcal F$ be a finite family of graphs of maximum degree at
most $d$, closed under induced subgraphs. Suppose that all $Z_H$ are
positive and, for every $H\in\mathcal F$ and $u\in V(H)$,
\begin{equation}\label{eq:second-order-assumptions}
 \frac{Z_H}{Z_{H-u}}\ge\frac18,
 \qquad
 \frac{\mathcal J_H(u)}{Z_H}\le2b_H(u)+\beta.
\end{equation}
For $z\ge R$ and either sign let
$G_{H,z}^{\pm}=(I\pm A_H/z)^{-1}$ on good signings.
Write $G_{H,z}=G_{H,z}^{-}$ when choosing one sign.
Define the finite maxima
\[
 F=\max_{H\in\mathcal F,\ u\in V(H)}
       \bigl\|(b_H(u)-(G_{H,\zeta})_{uu})_+\bigr\|_{L^2(\nu_H)},
 \qquad
 M=\max_{H\in\mathcal F,\ u\in V(H)}
       \mathbb E_{\nu_H}(G_{H,\zeta}^{\,2})_{uu}.
\]
An empty collection has maxima zero. Notice that the square in
$G_{H,\zeta}^{\,2}$ is a matrix square. Both maxima are finite:
$F\le B$ and $\|G_{H,\zeta}\|\le C_\zeta$, where
\[
 C_\zeta=\frac{\zeta}{\zeta-R}
 =\frac{1+k+\sqrt{1+k}}{k}\le\frac4k.
\]
The same deficit bound applies to the plus sign, because the Gibbs
measure is invariant under simultaneous reversal of every edge sign.

\begin{lemma}\label{lem:second-order-cavity}
Under these assumptions there is a constant $C_\tau$, depending only
on $\tau$, such that
\[
 F\le C_\tau\left(d^{-1/2}\beta^{-3/2}
                         +d^{-1}\beta^{-4}\right).
\]
If $K\in\mathcal F$ is extended by a new vertex with at most $d$
neighbors, let $s$ be its vector of independent incident signs and set
\[
 q_{R,\pm}=\frac1{R^2}s^T G_{K,R}^{\pm}s,
 \qquad q_* =\frac1{\zeta^2}\sum_{w\sim v}b_K(w).
\]
Then, under $\nu_K$ and independent incident signs,
\[
 \left(\mathbb E(q_*-q_{R,+})_+^2
             +\mathbb E(q_*-q_{R,-})_+^2\right)^{1/2}
 \le C_\tau\left(d^{-1/2}\beta^{-3/2}
                         +d^{-1}\beta^{-4}\right).
\]
\end{lemma}

\begin{proof}
We first bound the matrix-square moment by the deficit itself.
For each real $|\lambda|<R$, put $u=\lambda^2/R^2$ and
$\kappa=R/\zeta$. The even resolvent and even square-resolvent
kernels are
\[
 e_z(\lambda)=\frac1{1-\lambda^2/z^2},\qquad
 h_z(\lambda)=\frac{1+\lambda^2/z^2}{(1-\lambda^2/z^2)^2}.
\]
For $u>0$, direct division gives
\[
 \frac{h_\zeta(\lambda)-e_\zeta(\lambda)}
      {e_R(\lambda)-e_\zeta(\lambda)}
 =\frac{2\kappa^2(1-u)}{(1-\kappa^2)(1-\kappa^2u)}
 \le\frac2k.
\]
At $u=0$ both differences vanish. Applying spectral calculus and
then sign-reversal symmetry therefore gives
\begin{equation}\label{eq:ward-bootstrap}
 \mathbb E(G_{H,\zeta}^{\,2})_{uu}
 \le \mathbb E(G_{H,\zeta})_{uu}
   +\frac2k\left[\mathbb E(G_{H,R})_{uu}
                      -\mathbb E(G_{H,\zeta})_{uu}\right].
\end{equation}
The cofactor assumption and sign symmetry imply
$\mathbb E(G_{H,R})_{uu}\le b_H(u)+\beta/2$.
Also $\mathbb E(G_{H,\zeta})_{uu}\ge b_H(u)-F$.
Finally the even resolvent at $\zeta$ is pointwise bounded above
by the even resolvent at $R$, so sign symmetry gives
$\mathbb E(G_{H,\zeta})_{uu}\le B+\beta/2$.
Consequently
\begin{equation}\label{eq:second-moment-self-bound}
 M\le B+\frac\beta2+\frac\beta k+\frac{2F}{k}
   \le C_\tau\left(\beta^{-1}+\beta^{-2}F\right).
\end{equation}

Next fix $H,u$ and write $K=H-u$. With $s$ denoting the unnormalized
incident sign vector, put
\[
 D=(I-A_K/\zeta)^{-1},\qquad
 q=\zeta^{-2}s^TDs,\qquad
 \mu=\zeta^{-2}\sum_{w\sim u}D_{ww},\qquad
 q_0=\zeta^{-2}\sum_{w\sim u}b_K(w).
\]
Schur complements give $G_{uu}=1/(1-q)$ and
$b_H(u)=1/(1-q_0)$, where $G=G_{H,\zeta}$.
On the event $G_{uu}<b_H(u)$,
\[
 b_H(u)-G_{uu}=b_H(u)G_{uu}(q_0-q)
               \le B^2(q_0-q).
\]
The deterministic monotonicity $b_K(w)\le b_H(w)$ and the inverse
deletion formula
\[
 D_{ww}=G_{ww}-\frac{G_{wu}^2}{G_{uu}}
\]
give the pointwise estimate
\[
 q_0-q\le\frac1{\zeta^2}
             \sum_{w\sim u}(b_H(w)-G_{ww})_+
           +\frac{C_\zeta}{\zeta^2}+|q-\mu|.
\]
Here we used positivity of $G$ and
\[
 \sum_{w\sim u}\frac{G_{wu}^2}{G_{uu}}
 \le\frac{(G^2)_{uu}}{G_{uu}}\le\|G\|\le C_\zeta.
\]

Conditional on a fixed core, the exact Rademacher variance estimate
is
\[
 \mathbb E_{\mathrm{star}}(q-\mu)^2
 =\frac2{\zeta^4}\sum_{i\ne j\in N_H(u)}D_{ij}^2
 \le\frac2{\zeta^4}\sum_{i\in N_H(u)}(D^2)_{ii}.
\]
The full Gibbs density relative to $\nu_K$ and independent incident
signs is at most $8$, by the deletion-ratio assumption. Thus
\[
 \|q-\mu\|_{L^2(\nu_H)}
 \le\frac{\sqrt{16dM}}{\zeta^2}.
\]
Taking $L^2$ norms, applying Minkowski, and then maximizing over the
finite family gives
\[
 F\le B^2\left(pF+\frac{C_\zeta}{\zeta^2}
                         +\frac{4\sqrt{dM}}{\zeta^2}\right).
\]
Since $1-B^2p=\beta/a$ and $\zeta^{-2}=p/d\le1/(4d)$,
\begin{equation}\label{eq:cavity-l2-contraction}
 F\le\frac1{2\beta}
       \left(\frac{C_\zeta}{d}+4\sqrt{\frac Md}\right).
\end{equation}
Combining this with~\eqref{eq:second-moment-self-bound} yields
\[
 F\le C_\tau\left(
       d^{-1}\beta^{-3}
       +d^{-1/2}\beta^{-3/2}
       +d^{-1/2}\beta^{-2}\sqrt F\right).
\]
Young's inequality absorbs the last term into $F/2$ plus a constant
multiple of $d^{-1}\beta^{-4}$. Because $\beta\le1$, this proves
the stated bound for $F$.

For a fresh insertion, first consider the minus sign and put
$q_\zeta=\zeta^{-2}s^TG_{K,\zeta}s$.
For every $|\lambda|<R$, the function
$z\mapsto[z(z-\lambda)]^{-1}$ decreases for $z\ge R$.
Spectral calculus consequently gives $q_{R,-}\ge q_\zeta$.
The same statement for the plus sign follows by replacing $A_K$ by
$-A_K$. The fresh-star quadratic-form estimate, now without the
factor $8$ from density domination, gives
\[
 \|(q_*-q_{R,-})_+\|_2
 \le pF+\frac{\sqrt{2dM}}{\zeta^2}.
\]
By sign symmetry the plus-sign deficit has the same bound.
Use~\eqref{eq:second-moment-self-bound}, the established bound on $F$,
and $\sqrt F/(\sqrt d\,\beta)\le F/2+1/(2d\beta^2)$ to obtain the
claimed joint bound, after enlarging $C_\tau$.
\end{proof}


\section{Scalar insertion bounds}
\begin{lemma}[Coupling from squared lower-tail errors]
Let $q\in[0,1/2]$, $r=1-q$, and let $q_+,q_-\ge0$ be
random variables. Put
\[
 u=q_+-q,\qquad v=q_--q,\qquad
 \alpha=(r-u)_+(r-v)_+,
 \qquad D=(2r-u-v)\mathbf1_{\{u<r,\ v<r\}}.
\]
Suppose $M\ge0$ and $e\ge0$ satisfy
\[
 \mathbb E(u+v)\le M,
 \qquad \mathbb E(u_-^2+v_-^2)\le e^2.
\]
Set
\[
 C=e\sqrt{r(M+\sqrt2e)},\qquad f=\mathbb E\alpha.
\]
Then
\[
 f\ge r^2-rM-C,
 \qquad \mathbb E D-\frac{2f}{r}\le M+\frac{2C}{r}.
\]
In particular, if $f>0$ and
\begin{equation}\label{eq:second-order-scalar-budget}
 (1+\varepsilon r)M+(2/r+\varepsilon)C
 \le\varepsilon r^2,
\end{equation}
then $\mathbb E D/f\le2/r+\varepsilon$.
\end{lemma}

\begin{proof}
The appropriate cross term is capped:
\[
 L=u_-\min(v_+,r)+v_-\min(u_+,r).
\]
Pointwise on the whole probability space,
\begin{equation}\label{eq:second-order-pointwise}
 \alpha\ge r^2-r(u+v)-L,
 \qquad D-2\alpha/r\le u+v+2L/r.
\end{equation}
On $\{u<r,v<r\}$, expand
$\alpha=r^2-r(u+v)+uv$ and use $uv\ge-L$.
If, for example, $u\ge r$ and $v\ge0$, the right side of
the first inequality is nonpositive. If $u\ge r$ and $v<0$,
it is at most $rv_- - rv_-=0$. This proves the first inequality
also off the good event. For the second inequality, its left
side vanishes off the good event; there $u+v\ge r-q\ge0$,
because $u,v\ge-q$ and at least one is at least $r$.

Writing $P=\mathbb E(u_++v_+)$, one has
\[
 P=\mathbb E(u+v)+\mathbb E(u_-+v_-)
 \le M+\sqrt2e.
\]
Vector Cauchy--Schwarz and $\min(w,r)^2\le rw$ for $w\ge0$
give
\[
 \mathbb E L
 \le e\sqrt{\mathbb E\left[\min(u_+,r)^2+
                         \min(v_+,r)^2\right]}
 \le e\sqrt{rP}\le C.
\]
Averaging~\eqref{eq:second-order-pointwise} proves both estimates.
Finally~\eqref{eq:second-order-scalar-budget} is exactly
\[
 M+2C/r\le\varepsilon(r^2-rM-C)\le\varepsilon f.
\]
\end{proof}

\begin{proposition}[Explicit uniform parameter closure]
Let $0<\beta\le1/10$ and define
\[
 x=\frac{1-\beta^2}{4},\qquad
 a=\frac{1+\beta}{2},\qquad
 k=\frac{\beta^2}{1000},\qquad
 \zeta^2=\frac d x,\qquad R^2=\frac{\zeta^2}{1+k},
 \qquad x_R=\frac d{R^2}=(1+k)x.
\]
Suppose $a\le r\le1$, $q=1-r$, $0\le h\le x_R$, and set
\[
 \varepsilon=\beta,\qquad M=\beta h+2kq.
\]
If the hypotheses of the preceding lemma hold with
\[
 e\le\frac{\beta^{3/2}}{10},
\]
then
\[
 f>\frac15>\frac18,
 \qquad \frac{\mathbb E D}{f}\le\frac2r+\beta.
\]
Thus these assumptions close the deletion-ratio invariant
$f\ge c_0=1/8$ and the cofactor invariant
$\mathbb E D/f\le2/r+\varepsilon$.
\end{proposition}

\begin{proof}
Since $x_R=(1-\beta^2)(1+\beta^2/1000)/4\le1/4$ and $2q\le1$,
\[
 M\le m:=\beta x_R+k\le\frac{251}{1000}\beta.
\]
Also $\sqrt2e\le(45/1000)\beta$, so
\[
 C\le\frac1{10}\sqrt{\frac{296}{1000}}\,\beta^2
 <\frac{11}{200}\beta^2.
\]
Consequently
\[
 f\ge r^2-rM-C
 \ge\frac14-\frac{251}{1000}\beta-\frac{11}{200}\beta^2
 >\frac15.
\]

It remains to verify the scalar budget. The function
$r\mapsto\beta r^2-(1+\beta r)m$ is increasing for $r\ge a$,
since its derivative is $\beta(2r-m)>0$. Hence
\begin{align*}
 \beta r^2-(1+\beta r)M
 &\ge\beta a^2-(1+\beta a)(\beta x_R+k)\\
 &=\beta^2a(1-x_R)-k\bigl[1+\beta(a+x)\bigr]\\
 &\ge\left(\frac38-\frac{108}{100000}\right)\beta^2
 >\frac{37}{100}\beta^2.
\end{align*}
Here $a\ge1/2$, $x_R\le1/4$, and
$1+\beta(a+x)\le1+(1/10)(11/20+1/4)=108/100$.
On the other hand,
\[
 (2/r+\beta)C
 \le\frac{41}{10}\frac{11}{200}\beta^2
 =\frac{451}{2000}\beta^2
 <\frac{37}{100}\beta^2.
\]
This verifies~\eqref{eq:second-order-scalar-budget}.
\end{proof}


\section{Proof of the upper bound}
\begin{proof}[Proof of Theorem~\ref{thm:upper}]
Fix $\tau=1/1000$ and choose a sufficiently large absolute constant
$L$. Put $\beta=Ld^{-1/6}$ and take $d$ sufficiently large that
$\beta\le1/10$. Use the parameters of the preceding proposition;
in particular, $k=\tau\beta^2$ and $\zeta^2=4d/(1-\beta^2)$.
The cavity lemma implies that its joint fresh-star deficit $e$
satisfies
\[
 \frac e{\beta^{3/2}}
 \le C_\tau\left(L^{-3}+L^{-11/2}d^{-1/12}\right).
\]
Choosing $L$ and then $d$ sufficiently large makes this at most
$1/10$. These choices depend on no graph-order parameter.

Induct strongly on the number of vertices, proving simultaneously
for every graph of maximum degree at most $d$ that $Z_H>0$ and
\[
 Z_H/Z_{H-v}\ge1/8,
 \qquad \mathcal J_H(v)/Z_H\le2b_H(v)+\beta
 \quad(v\in V(H)).
\]
The empty graph is the base case. To insert $v$ into $H$, take
$K=H-v$ and apply Lemma~\ref{lem:second-order-cavity} to the finite
family of all induced subgraphs of $K$. Every hypothesis of that
lemma is already available by induction. The finite maxima within
its proof are over this known family; no property of $H$ is assumed.

Write $r=1/b_H(v)$ and $q=1-r$. By~\eqref{eq:path-recursion},
$q=\zeta^{-2}\sum_{w\sim v}b_K(w)$. Put $h=\deg_H(v)/R^2$.
The core cofactor bound and independent incident signs give
\begin{align*}
 \E_*(q_++q_-)
 &=R^{-2}\sum_{w\sim v}\mathcal J_K(w)/Z_K\\
 &\le2R^{-2}\sum_{w\sim v}b_K(w)+\beta h
   =2q+2kq+\beta h.
\end{align*}
The cavity lemma supplies the joint squared-deficit bound with
$e\le\beta^{3/2}/10$. The preceding proposition, together with
\eqref{eq:insertion-identities}, now gives $Z_H/Z_K>1/5$ and
$\mathcal J_H(v)/Z_H\le2/r+\beta$. This proves positivity and both
invariants, for every choice of $v$.

Positive $Z_H$ supplies a signing with norm strictly less than
\[
 R=\frac{2\sqrt d}{\sqrt{(1-\beta^2)(1+\beta^2/1000)}}
   =2\sqrt d+O(L^2d^{1/6}).
\]
Finally $2\sqrt d-2\sqrt{d-1}=O(d^{-1/2})$, which is absorbed in
the asserted additive bound.
\end{proof}
% Standalone section snippet.  The lower bound in this section has been
% independently checked algebraically and through the finite construction.
\section{An explicit positive excess in every degree}

Put $q=d-1\ge3$, and let $z_q$ be the unique positive solution of
\begin{equation}\label{eq:lower-root}
 qz(1+z+z^2+2z^3)=1+z+z^2+4z^4.
\end{equation}
Define
\[
 R_q=z_q^{-1/2}+qz_q^{1/2}.
\]
Then the worst optimal signing norm $\gamma_d$, whether defined over
finite simple graphs of maximum degree at most $d$ or over finite
simple $d$-regular graphs, satisfies
\begin{equation}\label{eq:explicit-lower}
 \gamma_d\ge R_q
 =2\sqrt q+(1+O(q^{-1}))q^{-11/2}.
\end{equation}
In particular, the exact Ramanujan bound fails in every degree
$d\ge4$, with an explicit positive quantitative excess.

To prove this, the polynomial obtained by subtracting the right side
of~\eqref{eq:lower-root} from its left side is
\[
 F_q(z)=-1+(q-1)z+(q-1)z^2+qz^3+(2q-4)z^4.
\]
It is strictly increasing for $z\ge0$, has $F_q(0)=-1$, and has
$F_q(1/q)=2q^{-3}-4q^{-4}>0$. Thus $0<z_q<1/q$.

Fix any threshold $R$ with $2\sqrt q<R<R_q$, and set
\[
 t=\frac{R-\sqrt{R^2-4q}}{2q},\qquad
 t_+=\frac{R+\sqrt{R^2-4q}}{2q}=\frac1{qt}.
\]
These are the two fixed points of $u\mapsto(R-qu)^{-1}$, and
$\sqrt{z_q}<t<q^{-1/2}$.

Construct a rooted seed from a triangle by attaching $q-2$ disjoint
complete rooted $q$-ary trees of height $h$ at its distinguished
vertex, and $q-1$ such trees at each of the other two vertices.
Each attachment is by one edge. The distinguished vertex has degree
$q$, all other nonterminal vertices have degree $q+1$, and terminal
vertices have degree one.

The seed has unsigned adjacency norm at most $2\sqrt q$. One direct
verification gives weight $1$ to its three triangle vertices and
weight $q^{-j/2}$ to vertices at distance $j\ge1$ down each attached
tree. At interior tree vertices the ratio $(Aw)_v/w_v$ is
$2\sqrt q$, and at leaves it is smaller. At the two nondistinguished
triangle vertices it is $2+(q-1)/\sqrt q\le2\sqrt q$, and at the
root it is smaller still. The positive-vector bound proves the claim.
The norm of every signing is no greater than the unsigned norm.
Consequently all seed resolvents at both $R$ and $-R$ exist.

The root response of a complete $q$-ary tree of height $h$ at $R$
converges increasingly to $t$. Signs on tree edges can be switched
away. Letting $h\to\infty$, eliminating these trees leaves effective
diagonal entries
\[
 a=R-(q-1)t=t^{-1}+t,
 \qquad b=R-(q-2)t=t^{-1}+2t.
\]
If the triangle sign product is $\eta\in\{-1,1\}$, the root
response is
\[
 g_\eta=\left(b-\frac2{a-\eta}\right)^{-1}.
\]
Changing the spectral side exchanges $\eta$ and $-\eta$. The mean
of the two limiting root responses is therefore independent of the
signing and equals
\begin{equation}\label{eq:triangle-mean}
 m(t)=\frac{g_++g_-}{2}
 =t\frac{1+t^2+t^4+2t^6}{1+t^2+t^4+4t^8}.
\end{equation}
By~\eqref{eq:lower-root},
\[
 m(t)>t_+\quad\Longleftrightarrow\quad F_q(t^2)>0.
\]
Hence for a sufficiently large finite height $h$, every signing of
the seed has mean two-sided root response greater than $t_+$.

Starting from this finite seed, repeatedly form a new rooted graph
by joining a new root to the roots of $q$ disjoint copies of the
previous graph. Suppose a signing of a graph built this way has norm
strictly less than $R$. All its principal resolvents are then
positive definite. If the child branches have mean two-sided root
response at least $u$, Schur complementation and the harmonic-mean
inequality give mean response at the new root at least
\[
 \frac1{R-qu}.
\]
Starting strictly above $t_+$, iterating this scalar map eventually
exceeds $R/(q+1)$, or a nonpositive required Schur denominator is
encountered earlier. Indeed, above $t_+$ the map is strictly larger
than its argument; an increasing bounded iteration would converge
to a fixed point, and neither fixed point is above $t_+$.
Joining $q+1$ copies of a sufficiently amplified branch at a final
central vertex is therefore impossible under the assumed norm bound:
the two central Schur inequalities would require the sum of the
child mean responses to be at most $R$.

This gives a finite connected simple graph $H$ of maximum degree
$d$ for which every signing has norm at least $R$. If regularity is
desired, take $d$ copies $H^{+,i}$ and $d$ copies $H^{-,i}$, with
$i\in\mathbb Z/d\mathbb Z$. For each original vertex $v$ of degree
$d-r_v$, add the edges
\[
 v^{+,i}v^{-,i+j},\qquad 0\le j<r_v.
\]
The resulting graph is simple and $d$-regular and contains each
original copy as an induced subgraph. It is connected: $H$ is
connected and has a leaf, whose deficiency $d-1\ge2$ supplies the
shifts $0$ and $1$ linking all copies. Its signing norm is at least
that of the restricted induced copy. Since $R\uparrow R_q$ is
arbitrary, this proves $\gamma_d\ge R_q$.

Finally, $F_q'(z)=q+O(1)$ uniformly near $z=1/q$, so
\[
 z_q=q^{-1}-2q^{-4}+O(q^{-5}).
\]
Writing $u_q=\sqrt{qz_q}=1-q^{-3}+O(q^{-4})$ gives
\[
 R_q-2\sqrt q
 =\sqrt q\,(u_q+u_q^{-1}-2)
 =\sqrt q\,\frac{(u_q-1)^2}{u_q}
 =(1+O(q^{-1}))q^{-11/2},
\]
as asserted.
% Stronger lower bound. Exact response map, domain, expansion, and finite
% approximation have been independently checked by two agents.
\section{A stronger lower bound from repeated triangle sources}

The preceding construction can be strengthened by repeating the
triangle source, with several tree levels between successive sources.
The resulting quantitative bound is
\begin{equation}\label{eq:buffered-main}
 \liminf_{q\to\infty}q^{7/2}
 \bigl(\gamma_{q+1}-2\sqrt q\bigr)\ge\frac14.
\end{equation}
We give the finite construction and the response estimates explicitly.

Fix $0<c<1/4$, put $s=\sqrt c$, and choose a fixed integer $L\ge2$
such that
\begin{equation}\label{eq:buffer-choice}
 s<\frac{L-1}{2(L+1)}.
\end{equation}
For all sufficiently large integers $q$, we will construct a finite
simple graph of maximum degree $q+1$ for which no signing has norm
strictly less than
\[
 R=2\sqrt q+cq^{-7/2}.
\]
Use normalized root responses $g_\pm=\sqrt q\,
[(RI\mp A_\sigma)^{-1}]_{oo}$ and put
\[
 z=R/\sqrt q=2+cq^{-4},\qquad
 \phi(t)=\frac1{z-t},\qquad
 \tau=\frac{z-\sqrt{z^2-4}}2,\qquad
 y_+=\frac{z+\sqrt{z^2-4}}2.
\]
Thus $\tau$ and $y_+$ are the two fixed points of $\phi$.
All response identities below are used under the assumption that the
signed graph in question has norm less than $R$. Its principal
resolvents are then positive definite, so every stated denominator
has the indicated positive sign.

\subsection*{A tree buffer retains information about both responses}

A $q$-join adds a root adjacent to the roots of $q$ disjoint rooted
graphs. At the new root, each normalized response is $\phi$ applied
to the average of the corresponding child responses. Let
\[
 F_L=\phi^{-L}=(t\mapsto z-t^{-1})^{\circ L}.
\]
This is increasing and concave on the interval
$(\phi^{\circ L}(0),\infty)$. Genuine responses after $L$ joins
belong to this interval: at every join the average child response is
positive. If every signing of the initial rooted graph satisfies
$(g_++g_-)/2\ge y$, then the two responses $\xi_\pm$ after $L$
successive $q$-join levels satisfy
\begin{equation}\label{eq:buffer-constraint}
 F_L(\xi_+)+F_L(\xi_-)\ge2y.
\end{equation}
Indeed, at one join $F_1$ recovers the average child response; at
subsequent joins use concavity of $F_j$ and Jensen's inequality.
In particular, averaging any number of such buffered response pairs
preserves~\eqref{eq:buffer-constraint}.

\subsection*{The repeated source and its exact scalar lower bound}

Given a rooted graph $X$ whose root has degree $q$, form a new rooted
graph as follows. Take a triangle and distinguish one of its
vertices as the new root. Attach $q-2$ disjoint copies of an
$L$-level $q$-join buffer of $X$ to this root. At each of the two
other triangle vertices attach $q-1$ disjoint complete $q$-ary trees.
For the moment use the limiting normalized tree response $\tau$;
the trees will be made finite below. All vertices have degree at
most $q+1$, and the distinguished root again has degree $q$.

Write
\[
 b=1-\frac2q,\qquad
 k_\pm=\frac{2/q}{z-(1-1/q)\tau\mp q^{-1/2}},
 \qquad a_\pm=z-k_\pm.
\]
After eliminating the pure tree branches and the two other triangle
vertices, the two normalized root responses are
$[a_\pm-b\bar\xi_\pm]^{-1}$, up to exchanging $+$ and $-$ according
to the triangle sign product. Here $\bar\xi_\pm$ are the averages
of the buffered child responses.
By~\eqref{eq:buffer-constraint}, setting
$t_\pm=F_L(\bar\xi_\pm)$ gives $t_++t_-\ge2y$. Therefore the new
mean response is bounded below by the minimum of
\[
 \frac{\chi_+(t_+)+\chi_-(t_-)}2,
 \qquad
 \chi_\pm(t)=\frac1{a_\pm-b\phi^{\circ L}(t)},
 \qquad t_++t_-\ge2y.
\]

Here is an exact formula for this minimum. Write
\[
 \begin{pmatrix}A&B\\ C&D\end{pmatrix}
 =\begin{pmatrix}0&1\\-1&z\end{pmatrix}^{L},
 \qquad \phi^{\circ L}(t)=\frac{At+B}{Ct+D}.
\]
The determinant is one. Define
\[
 E_\pm=a_\pm C-bA,\qquad H_\pm=a_\pm D-bB,
 \qquad E_0=(E_++E_-)/2,\quad H_0=(H_++H_-)/2,
 \quad\Delta=a_+-a_-.
\]
For fixed $L$ and large $q$, $C<0$ and $E_\pm<0$. On the lower
positive-denominator branch,
\[
 \chi_\pm(t)=\frac{Ct+D}{E_\pm t+H_\pm},\qquad
 \chi_\pm'(t)=\frac b{(E_\pm t+H_\pm)^2}>0,
 \qquad \chi_\pm''(t)>0.
\]
Their poles are $1+1/(L+1)+o(1)$, which lie below the common
numerator zero $1+1/L+o(1)$. The functions are therefore positive
and convex on the entire intervals below their respective poles.
Thus allowing this entire branch, instead of just the smaller set
of attainable responses, gives a valid lower bound. Monotonicity
reduces the minimization to $t_++t_-=2y$. Equal derivatives then
give the unique minimizer $t_\pm=y\pm w$, where
\[
 w=-\frac{\Delta(Cy+D)}{2E_0}.
\]
For $y$ near $1$ this minimizer lies strictly below both poles.
Substitution yields the exact lower-bound map
\begin{equation}\label{eq:buffer-map}
 M_L(y)=
 \frac{Cy+D}
 {E_0y+H_0-\displaystyle\frac{\Delta^2 C(Cy+D)}{4E_0}}.
\end{equation}
The terms involving $\Delta^2$ cancel in the determinant, so
\begin{equation}\label{eq:buffer-map-derivative}
 M_L'(y)=
 \frac b
 {\left(E_0y+H_0-\displaystyle\frac{\Delta^2 C(Cy+D)}{4E_0}\right)^2}.
\end{equation}
Consequently the repeated source takes a uniform lower bound $y$
on mean response to the uniform lower bound $M_L(y)$.

\subsection*{Expansion at the upper tree fixed point}

Set $e=q^{-1/2}$, so $z=2+s^2e^8$. With $L$ and $s$ fixed,
\[
 \tau=1-se^4+O(e^8),\qquad
 y_+=1+se^4+O(e^8),
\]
and direct expansion gives
\[
 k_\pm=2e^2\pm2e^3\mp2e^5-2(1+s)e^6+O(e^7).
\]
The optimizer in~\eqref{eq:buffer-map} has
$w=-2e^3/(L+1)+O_L(e^5)$. Since $\phi(y_+)=y_+$,
uniformly for $w=O_L(e^3)$,
\[
 \phi^{\circ L}(y_+\pm w)
 =y_+\pm w+Lw^2+O_L(e^7).
\]
Substituting into the two root responses and averaging gives
\[
 \frac{\chi_+(y_++w)+\chi_-(y_+-w)}2-y_+
 =Lw^2+(w+2e^3)^2-(2+4s)e^6+O_L(e^7).
\]
Minimizing the displayed quadratic, equivalently using the exact
optimizer, proves
\begin{equation}\label{eq:buffer-positive-drift}
 M_L(y_+)-y_+
 =\left(\frac{4L}{L+1}-2-4s\right)q^{-3}
   +O_L(q^{-7/2}).
\end{equation}
The leading coefficient is strictly positive by
\eqref{eq:buffer-choice}.
Moreover, uniformly for $y\in[\tau-O(q^{-3}),y_+]$,
\begin{equation}\label{eq:buffer-contraction}
 M_L'(y)=1-\frac2q+O_L(q^{-2})<1
\end{equation}
for sufficiently large $q$. To see the expansion directly in
\eqref{eq:buffer-map-derivative}, its denominator before squaring is
$1+O_L(q^{-2})$ on this interval. This follows from
$A=-(L-1)+O_L(q^{-4})$, $B=L+O_L(q^{-4})$,
$C=-L+O_L(q^{-4})$, $D=L+1+O_L(q^{-4})$,
$a_++a_-=4-4/q+O(q^{-3})$, and $\Delta=O(q^{-3/2})$.
It follows that $M_L(y)-y$ is decreasing on the stated interval.
Together with~\eqref{eq:buffer-positive-drift}, this gives a uniform
strict positive drift all the way from the lower tree fixed point
to the upper one.

\subsection*{All trees and all iterations can be finite}

Fix a sufficiently large $q$. Choose $\beta>0$, independent of $q$
once $L,s$ are fixed, so that
\[
 M_L(y)-y\ge\beta q^{-3}
 \quad\hbox{for }y\in I:=[\tau-\beta q^{-3}/8,y_+].
\]
The normalized response $\tau_h$ of a finite complete $q$-ary tree
of height $h$ increases to $\tau$. Choose $h$ so large that
\[
 \tau_h\in I,\qquad
 \sup_{y\in I}|M_{L,h}(y)-M_L(y)|<\beta q^{-3}/4,
\]
where $M_{L,h}$ is~\eqref{eq:buffer-map} with $\tau$ replaced by
$\tau_h$ in $k_\pm$. This is possible by continuity of its rational
coefficients and because all denominators stay bounded away from
zero on the compact interval $I$. The same lower-branch
minimization remains valid for these finite tails.

Start $X_0$ as a finite complete $q$-ary tree with normalized root
response $\tau_h$. Repeatedly apply the triangle operation described
above, always using pure tails of this same finite height. If a
resulting signing has norm less than $R$, the preceding estimates
give lower bounds $y_0=\tau_h$, $y_{j+1}=M_{L,h}(y_j)$ on its mean
root responses. While $y_j\le y_+$ these increase by at least
$3\beta q^{-3}/4$. Hence after finitely many steps (in fact
$O_L(q)$ steps) the lower bound exceeds $y_+$.

Finally, ordinary $q$-joins amplify any mean normalized response
strictly above $y_+$ by the map $\phi$. After finitely many such
joins it exceeds $qz/(q+1)$. Attaching $q+1$ resulting branches to
a central vertex contradicts the two positive Schur-complement
inequalities at that vertex. Thus the final finite core has no
signing of norm less than $R$. Its maximum degree is $q+1$.
The finite induced regular completion used in the preceding section
produces a connected simple $(q+1)$-regular graph with the same
obstruction.

We have therefore proved, for each fixed $c<1/4$ and all sufficiently
large $q$, that
\[
 \gamma_{q+1}\ge2\sqrt q+cq^{-7/2}.
\]
Letting $c\uparrow1/4$ proves~\eqref{eq:buffered-main}.

\section{What remains unresolved}
The upper bound proved here does not give a bounded additive error,
much less identify an asymptotic equivalent for
$\gamma_d-2\sqrt{d-1}$. The exponent $1/6$ comes from the balance
\[
 F=O(d^{-1/2}\beta^{-3/2}+d^{-1}\beta^{-4}),\qquad
 \text{insertion error}=O(F\sqrt\beta),\qquad
 \text{available margin}\asymp\beta^2.
\]
Thus this method takes $\beta\asymp d^{-1/6}$ and spectral excess
$\beta^2\sqrt d\asymp d^{1/6}$. This calculation describes the
present proof's limitations; it is not a lower-bound obstruction
at that scale. The argument is existential and does not establish
a polynomial-time algorithm.

\begin{thebibliography}{9}
\bibitem{JSS} A. Jadbabaie, A. Saberi, and S. Sra,
\emph{Kadison--Singer partitions and Bilu--Linial graph signings in
polynomial time}, arXiv:2609.23855v1 (2026).
\url{https://arxiv.org/abs/2609.23855}.
\bibitem{Xu} Z. Xu, \emph{A 3-regular counterexample to the Bilu--Linial
signing conjecture}, arXiv:2609.15591v3 (2026).
\url{https://arxiv.org/abs/2609.15591}.
\end{thebibliography}
\end{document}
